\documentclass[12pt, oneside, a4paper]{ctexart}
\usepackage{amsmath}
\usepackage{amssymb}
\usepackage{ctex}
\usepackage{enumerate}
\usepackage[{top=2.54cm, bottom=2.54cm, left=3.18cm, right=3.18cm}]{geometry}
\usepackage{graphicx}
\usepackage[breaklinks, colorlinks=false, hidelinks]{hyperref}
\usepackage{mathptmx}
\usepackage{xcolor}
\usepackage{xeCJK}

\ctexset{
    section={
        name={,、},
        number={\chinese{section}},
        format=\heiti\raggedright\zihao{3}\bfseries,
        aftername={\hspace{0.5em}},
        beforeskip=1ex,
        afterskip=1ex,
    },
    subsection={
        name={,.},
        number={\arabic{subsection}},
        format=\songti\raggedright\zihao{-4}\bfseries,
        aftername={\hspace{0.5em}},
        beforeskip=0ex,
        afterskip=0ex,
    },
}
\pagestyle{plain}
\renewcommand{\figurename}{\small{图}}
\setmainfont{Times New Roman}
\setCJKmainfont[AutoFakeBold={2}]{SimSun}
\setCJKfamilyfont{zh}{SimSun}
\setCJKfamilyfont{ja}{BIZ UDMincho}
\xeCJKDeclareCharClass{CJK}{
  "24EA,
  "2460->"2473,
  "3251->"32BF,
  "24FF,
  "2776->"277F,
  "24EB->"24F4
}


%----------分割线----------%


\begin{document}


\thispagestyle{empty}
\begin{center}{
    \fontsize{26pt}{26pt}\selectfont
    \vspace*{2.75\baselineskip}
    \includegraphics[width=0.6\textwidth]{pics/tongji-university-logo-1024px.png}\\
    [\baselineskip]
    \heiti{轨道交通综合实验报告（一）}\\
    \fontsize{18pt}{24pt}\selectfont
    \vspace*{1.5\baselineskip}
    \songti{\begin{tabular}{l}
            {\heiti{}实验日期：}【待填写实际日期】    \\
            {\heiti{}实验名称：}振动测试技术实验     \\
            {\heiti{}实验人员：}肖哲夫~~2352670 \\
            {\heiti{}实验地点：}【待填写实际地点】
        \end{tabular}
    }
    }
\end{center}
\newpage{}
\setcounter{page}{1}

\section{\heiti{实验目的}}
\par{}
掌握加速度传感器的使用；掌握振动测量技术；掌握振动分析技术。


\section{\heiti{实验设备}}
\par{}
等强度梁、加速度传感器、电荷放大器、计算机、数据采集卡和不同质量的砝码。
\par{}
数据采集程序准备阶段使用信号发生器、USB-6008数据采集卡及安装了LabVIEW和相应驱动程序的计算机。%各设备的型号与实际参数见实验报告部分。

\section{\heiti{实验步骤}}
\subsection{\heiti{加速度传感器介绍}}
\par{}
由指导教师介绍加速度传感器的技术指标，包括量程、测量方向和响应频率；由指导教师介绍加速度传感器的标定系数和标定方法。
\subsection{\heiti{在等强度梁上粘贴加速度传感器}}
\par{}
根据测量方向，使用磁铁或502胶水将加速度传感器固接在等强度上。
\subsection{\heiti{测量仪器的使用}}
\par{}
由指导教师介绍加速度传感器、电荷放大器的使用方法。
\subsection{\heiti{计算机测试系统的使用}}
\par{}
以往的数据采集与信号处理的程序；标定系数的设定。
\subsection{\heiti{振动加速度的测试}}
\par{}
连接加速度传感器、电荷放大器和计算机测试系统；对振动测试系统进行调试、标定；对等强度梁进行激振，记录加速度振动波形，对加速度振动波形进行频谱分析；给等强度梁上增加砝码，重新测量振动频率，记录波形；不同砝码重量的重复实验三次，记录波形。



\section{实验报告}
\subsection{\heiti{加速度传感器的使用方法}}
\par{}
【待根据实际设备填写】%本实验使用加速度传感器测量等强度梁的振动加速度。传感器应沿预定测量方向固定在梁上，并通过电荷放大器接入数据采集系统。传感器型号、量程、测量方向、响应频率、安装位置及固定方式分别为：\textcolor{gray}{【待根据实际设备填写】}。
\subsection{\heiti{测试系统的标定系数}}
\par{}
【待填写】
%记录传感器标定系数、电荷放大器设置及采集程序中使用的换算系数。若整个测试系统以$K_a$（单位为$g/\mathrm{V}$）表示电压至加速度的换算，则处理后的加速度可写为
%\begin{equation}
%    a(t)=K_a\bigl[u(t)-u_0\bigr],
%\end{equation}
%其中$u(t)$为采集电压，$u_0$为零位电压。实际使用的标定系数、单位及设置位置为：\textcolor{gray}{【待填写】}。
\subsection{\heiti{实验原始数据与加速度振动波形}}
\par{}
【待填写】
%记录各工况的砝码质量、采样设置和原始波形文件。每种砝码质量下的三次重复试验应分别对应一份波形及其文件名，不得用处理后的结果代替原始记录。
\par{}
\begin{center}
    \renewcommand{\arraystretch}{1.3}
    \small
    \begin{tabular}{|c|c|c|c|c|}
        \hline
        工况    & 附加砝码质量 & 第一次波形文件 & 第二次波形文件 & 第三次波形文件 \\
        \hline
        基准工况  & 【待填写】  & 【待填写】   & 【待填写】   & 【待填写】   \\
        \hline
        加砝码工况 & 【待填写】  & 【待填写】   & 【待填写】   & 【待填写】   \\
        \hline
    \end{tabular}
\end{center}
\par{}
\textcolor{gray}{【待补充：按照实际砝码质量增减表格行，并逐一插入所有激振波形；图题注明工况和试验次数。】}
\subsection{\heiti{数据处理结果与振动频率分析}}
\par{}
【待填写】
%对每次采集的加速度波形进行频谱分析，记录主要频率峰值。若对同一工况进行三次重复测量，其平均频率为
%\begin{equation}
%    \bar f=\frac{f_1+f_2+f_3}{3}.
%\end{equation}
%频谱分析方法、频谱图和测量结果为：\textcolor{gray}{【待根据原始数据补充】}。
\par{}
\begin{center}
    \renewcommand{\arraystretch}{1.3}
    \small
    \begin{tabular}{|c|c|c|c|c|c|}
        \hline
        附加砝码质量 & $f_1$ / Hz & $f_2$ / Hz & $f_3$ / Hz & $\bar f$ / Hz & 备注    \\
        \hline
        【待填写】  & 【待填写】      & 【待填写】      & 【待填写】      & 【待填写】         & 【待填写】 \\
        \hline
        【待填写】  & 【待填写】      & 【待填写】      & 【待填写】      & 【待填写】         & 【待填写】 \\
        \hline
    \end{tabular}
\end{center}
\subsection{\heiti{实验结论}}
\par{}
\textcolor{gray}{【待根据实测波形、频率与重复测量结果填写结论，说明不同附加质量下的频率变化及可能的误差来源。】}
\subsection{\heiti{思考题}}
\par\noindent{}
\heiti{(1)~砝码质量增加1倍时，等强度梁的振动频率增加多少倍？分析测量结果与理论结论出现差异的原因。}
\par{}
{\setlength{\parindent}{0pt}\heiti{}答：}
\par{}
%单自由度近似下，振动频率满足$f=(2\pi)^{-1}\sqrt{k/m_{\mathrm{eq}}}$。若仅砝码质量由$m$变为$2m$，而梁和传感器等形成的其余等效质量为$m_0$，则频率比为
%\begin{equation}
%    \frac{f_2}{f_1}=\sqrt{\frac{m_0+m}{m_0+2m}}.
%\end{equation}
%当其余等效质量相对砝码质量可忽略时，该比值趋近于$1/\sqrt{2}$，即频率降低而非增加。实测频率比和偏差原因需结合实验结果分析：\textcolor{gray}{【待填写】}。
\par\noindent{}
\heiti{(2)~已知砝码质量时，如何通过系统振动频率推算等强度梁的刚度？}
\par{}
{\setlength{\parindent}{0pt}\heiti{}答：}
\par{}
%若已知振动系统的等效质量$m_{\mathrm{eq}}$，由单自由度模型可得等效刚度
%\begin{equation}
%    k=(2\pi f)^2m_{\mathrm{eq}}.
%\end{equation}
%若梁和传感器的等效质量未知，可利用两种已知附加质量$m_1$、$m_2$及其对应频率$f_1$、$f_2$消去固定等效质量：
%\begin{equation}
%    k=\frac{4\pi^2(m_2-m_1)}{1/f_2^2-1/f_1^2}.
%\end{equation}
%根据实际测量数据计算的刚度为：\textcolor{gray}{【待填写】}。
\subsection{\heiti{程序源代码}}
\par{}
\textcolor{gray}{【待附实际使用的LabVIEW VI程序及程序框图，注明采集、标定、波形文件记录和频谱分析部分。】}
\subsection{\heiti{附上含实验人员和日期的现场照片。}}
\par{}
\textcolor{gray}{【待插入现场照片；图题注明实验人员、拍摄日期和实验环节。】}

\end{document}
